\documentclass[10pt,a4paper]{article}
\usepackage[utf8]{inputenc}
\usepackage[T1]{fontenc}
\usepackage{lmodern,amsmath,amssymb,booktabs,microtype,hyperref,needspace}
\usepackage[margin=22mm]{geometry}
\hypersetup{colorlinks=true,urlcolor=blue,linkcolor=black,pdftitle={Paper Zero - revised manuscript 2026-10-28},pdfauthor={Graziano Guiducci}}
\setlength{\parindent}{0pt}
\setlength{\parskip}{5pt}
\setlength{\emergencystretch}{2em}
\raggedbottom
\begin{document}
\begin{center}{\LARGE\bfseries L'incompletezza generativa}\end{center}

\textbf{Graziano Guiducci \textperiodcentered{}  Manoscritto rivisto \textperiodcentered{}  28 ottobre 2026}
\par

\emph{Una ricorsione. Una direzione invariante. Uno stato che prosegue.}
\par

La sorgente D-ND descrive il movimento della possibilità attraverso un riferimento indeterminato e una relazione dipolare che contiene i suoi poli. L'osservatore è l'ente risultante dell'allineamento; osservatore, allineamento e risultante si descrivono nello stesso movimento. La matematica è compresa in questo modello logico generale. Paper Zero ne presenta una particolare rappresentazione matematica: una ricorsione razionale con una direzione proiettiva attrattiva irrazionale.
\par

Il calcolo rende osservabile una relazione precisa: una direzione può restare invariante mentre lo stato che la porta continua a trasformarsi.
\par

\Needspace{9\baselineskip}\section*{La mappa}

\[
f(x)=1+\frac{1}{x}.
\]

Per un valore iniziale positivo, ogni iterata affine è definita e la successione converge al rapporto aureo. La tabella comprende anche un valore iniziale negativo la cui specifica orbita rimane definita. I valori sono arrotondati a tre cifre decimali.
\par

\begin{center}
\begin{tabular}{lrrr}
\toprule
Passo & $x_0=7$ & $x_0=0.1$ & $x_0=-3$ \\
\midrule
0 & 7.000 & 0.100 & -3.000 \\
1 & 1.143 & 11.000 & 0.667 \\
2 & 1.875 & 1.091 & 2.500 \\
3 & 1.533 & 1.917 & 1.400 \\
4 & 1.652 & 1.522 & 1.714 \\
5 & 1.605 & 1.657 & 1.583 \\
6 & 1.623 & 1.603 & 1.632 \\
7 & 1.616 & 1.624 & 1.613 \\
8 & 1.619 & 1.616 & 1.620 \\
\bottomrule
\end{tabular}
\end{center}

I due punti fissi reali sono
\par

\[
\varphi=\frac{1+\sqrt5}{2},\qquad \psi=\frac{1-\sqrt5}{2}=-\frac1\varphi.
\]

Sulla retta reale proiettiva, la trasformazione è definita in ogni punto e converge alla direzione rappresentata da $\varphi$ da ogni direzione iniziale, eccetto il repellente fisso $\psi$. Nella carta affine, lo zero e le sue preimmagini finite incontrano una divisione per zero: per esempio, $-1\mapsto0$ non ha un valore affine successivo. L'orbita proiettiva prosegue attraverso il punto all'infinito.
\par

\Needspace{9\baselineskip}\section*{La matrice e la convergenza}

\[
M=\begin{pmatrix}1&1\\1&0\end{pmatrix},\qquad \operatorname{tr}M=1,\qquad \det M=-1.
\]

Il polinomio caratteristico è $\lambda^2-\lambda-1$. I suoi autovalori sono $\varphi$ e $\psi$. Entrambi sono irrazionali algebrici quadratici.
\par

Siano $r=(\varphi,1)^T$ e $s=(\psi,1)^T$. Ogni vettore non nullo si scrive come $v=ar+bs$. Per $a\ne0$,
\par

\[
[M^n v]=\left[r+\frac ba\left(\frac\psi\varphi\right)^n s\right]\longrightarrow[r],\qquad \left|\frac\psi\varphi\right|=\varphi^{-2}<1.
\]

Le parentesi quadre indicano una direzione proiettiva: multipli scalari non nulli rappresentano lo stesso punto. Se $a=0$, la direzione è $[s]$ e rimane fissa. Questo dimostra la convergenza proiettiva enunciata, compresa la sua eccezione.
\par

Le iterate affini razionali finite rimangono razionali quando sono definite. Il loro limite attrattivo è irrazionale. Il titolo nomina questa relazione fra il campo razionale delle iterate finite e la direzione invariante nella sua estensione quadratica $\mathbb Q(\sqrt5)$. Questo è l'enunciato algebrico qui dimostrato. I teoremi di incompletezza di Gödel riguardano la dimostrabilità nelle teorie formali sotto le proprie ipotesi.
\par

\Needspace{9\baselineskip}\section*{Orientazione e crescita}

Il determinante esprime una proprietà geometrica esatta: per qualsiasi coppia di vettori, l'area orientata che essi generano cambia segno sotto $M$, mentre il suo valore assoluto è conservato.
\par

Crescita e periodicità dipendono anche dalla struttura degli autovalori. Per confronto,
\par

\[
B=\begin{pmatrix}0&1\\1&0\end{pmatrix},\quad B^2=I,\quad \det B=-1;\qquad C=\begin{pmatrix}2&1\\1&1\end{pmatrix},\quad \det C=1.
\]

$B$ è periodica; $C$ ha autovalore espansivo $(3+\sqrt5)/2$. Nella coppia reciproca scelta, cambiare la regola in $g(x)=1-1/x$ dà
\par

\[
N=\begin{pmatrix}1&-1\\1&0\end{pmatrix},\qquad N^3=-I.
\]

La sua trasformazione proiettiva ha periodo tre. Questo confronto distingue queste due regole. Il solo segno del determinante non classifica la crescita generativa.
\par

\Needspace{9\baselineskip}\section*{Perché compare il rapporto aureo}

L'equazione del punto fisso positivo $x=1+1/x$ dà $x^2-x-1=0$ e quindi $x=\varphi$. La sua frazione continua è $[1;1,1,1,\ldots]$.
\par

Si può enunciare una minimalità precisa: fra le matrici intere non negative primitive $2\times2$ con determinante $-1$, il raggio spettrale minimo è $\varphi$. Infatti la loro traccia intera $t$ è almeno uno e il raggio spettrale è $(t+\sqrt{t^2+4})/2$. Per $t=1$, gli elementi diagonali sono zero e uno e quelli fuori diagonale sono entrambi uno. I due rappresentanti sono $M$ e $PMP$, dove $P=B$ scambia le coordinate.
\par

L'entropia appartiene a un sistema dinamico specificato. Con $M$ come matrice di adiacenza, lo shift golden mean ha entropia topologica $\log\varphi$ (logaritmo naturale). La sostituzione Fibonacci con la stessa matrice di incidenza ha entropia zero. La mappa proiettiva qui studiata è un omeomorfismo della circonferenza e ha anch'essa entropia zero. Sono impieghi diversi della stessa matrice.
\par

\Needspace{9\baselineskip}\section*{Una direzione fissa e uno stato che prosegue}

Sulla direzione attrattiva,
\par

\[
Mr=\varphi r,\qquad M^nr=\varphi^nr,\qquad [M^nr]=[r].
\]

Il vettore cambia scala a ogni passo; la sua direzione proiettiva resta fissa. Questa relazione era già sviluppata nel manoscritto D-ND sulla distinzione relazionale e la continuità della risultante, Appendice B, 20 ottobre 2026.
\par

In questa rappresentazione, l'espressione $R+1=R$ si può leggere come un'ulteriore applicazione della trasformazione proiettiva che lascia invariata la sua direzione fissa: $T([r])=[r]$. È una lettura matematica della risultante D-ND. Il significato sorgente di risultante, riferimento, movimento e osservatore rimane raggiungibile oltre questa rappresentazione scelta.
\par

\Needspace{9\baselineskip}\section*{Le due coordinate}

Si scrive lo stato come coppia di coordinate in relazione $(D,ND)$. Allora
\par

\[
(D,ND)\longmapsto(D+ND,D),\qquad x=\frac D{ND},\qquad x'=\frac{D+ND}{D}=1+\frac1x.
\]

Entrambe le coordinate sono già presenti nell'ingresso e vengono aggiornate in una sola operazione. I loro nomi forniscono l'interpretazione scelta per questa costruzione; la relazione dipolare che contiene i poli appartiene alla sorgente D-ND più ampia. Il determinante descrive l'area, anziché la conservazione del prodotto delle coordinate $D\cdot ND$.
\par

\Needspace{9\baselineskip}\section*{Il primo passo numerico}

Da $x_0=7$, la prima iterata è esattamente $8/7$ e la quarta è $38/23$. Rispetto all'errore assoluto iniziale $|7-\varphi|$, l'errore diminuisce di circa \textbf{91.170943\%} dopo un passo e \textbf{99.365661\%} dopo quattro passi. I consueti valori arrotondati, \textbf{91.2\%} e \textbf{99.4\%}, esprimono questo esempio numerico. Le percentuali misurano la riduzione dell'errore assoluto.
\par

\Needspace{9\baselineskip}\section*{La relazione resa visibile}

Una matrice intera possiede un'autodirezione attrattiva irrazionale. L'iterazione proiettiva razionale si avvicina a quella direzione nel dominio dichiarato; un autovettore su di essa continua a cambiare scala. Paper Zero rende esplicita questa relazione fra invarianza e continuazione nella ricerca D-ND più ampia.
\par

La prova algebrica è esposta sopra. L'iterazione ne rende osservabile la conseguenza.
\par

\Needspace{9\baselineskip}\section*{Fonti ed edizione}

Il modello D-ND e le sue osservazioni sorgente sono attribuiti a Graziano Guiducci. Il chiarimento autoriale del 27 ottobre 2026 colloca la matematica nel modello logico generale. Questo manoscritto rivisto usa matematica classica; esempi matematici e interpretazioni sono distinti da quelle formulazioni autoriali.
\par

Per la distinzione fra entropie: H. Bruin, \emph{Ergodic Theory I}, lezione 18, Università di Vienna; M. Baake, T. Spindeler e N. Strungaru, \emph{Diffraction of compatible random substitutions in one dimension}, Indagationes Mathematicae 29 (2018), 1031--1071, DOI: 10.1016/j.indag.2018.05.008. L'enunciato sull'entropia proiettiva segue anche direttamente dall'argomento sulle partizioni della circonferenza riportato nelle note di revisione allegate.
\par

Edizione storica: marzo 2026, DOI \textbf{10.5281/zenodo.18902950}, archivio della release v2.0.0 del repository. Quel DOI identifica il deposito storico; questo manoscritto rivisto conserva una propria identità datata.
\par

\Needspace{9\baselineskip}\section*{Informazioni sul deposito}

\textbf{Licenza del manoscritto:} Copyright 2026 Graziano Guiducci. \href{https://creativecommons.org/licenses/by/4.0/}{CC BY 4.0}. La licenza riguarda questo manoscritto e i corrispondenti sorgenti Markdown e LaTeX.
\par

\textbf{DOI del manoscritto:} \href{https://doi.org/10.5281/zenodo.23735742}{10.5281/zenodo.23735742}.
\par
\end{document}
